\documentclass[10pt,a4paper]{amsart}
\usepackage{amsmath,amssymb,amsthm,graphicx}
% Retained minimal preamble for the Stacks Project chapter varieties.tex.
% Verbatim excerpt of the chapter preamble: Xy-pic setup (native preamble lines 14-18), the native theorem
% declarations with the mathematical macros (native preamble lines 154-217), and the project own no-op
% definitions of the reference/slogan/history display environments (native wrapper.tex lines 1-3) that the
% counted statement uses for its slogan. Copied mechanically; nothing added, removed or redefined.
% For commutative diagrams we use Xy-pic
\usepackage[all]{xy}

% We use 2cell for 2-commutative diagrams.
\xyoption{2cell}
% Theorem environments.
%
\theoremstyle{plain}
\newtheorem{theorem}[subsection]{Theorem}
\newtheorem{proposition}[subsection]{Proposition}
\newtheorem{lemma}[subsection]{Lemma}

\theoremstyle{definition}
\newtheorem{definition}[subsection]{Definition}
\newtheorem{example}[subsection]{Example}
\newtheorem{exercise}[subsection]{Exercise}
\newtheorem{situation}[subsection]{Situation}

\theoremstyle{remark}
\newtheorem{remark}[subsection]{Remark}
\newtheorem{remarks}[subsection]{Remarks}

\numberwithin{equation}{subsection}

% Macros
%
\def\lim{\mathop{\mathrm{lim}}\nolimits}
\def\colim{\mathop{\mathrm{colim}}\nolimits}
\def\Spec{\mathop{\mathrm{Spec}}}
\def\Hom{\mathop{\mathrm{Hom}}\nolimits}
\def\Ext{\mathop{\mathrm{Ext}}\nolimits}
\def\SheafHom{\mathop{\mathcal{H}\!\mathit{om}}\nolimits}
\def\SheafExt{\mathop{\mathcal{E}\!\mathit{xt}}\nolimits}
\def\Sch{\mathit{Sch}}
\def\Mor{\mathop{\mathrm{Mor}}\nolimits}
\def\Ob{\mathop{\mathrm{Ob}}\nolimits}
\def\Sh{\mathop{\mathit{Sh}}\nolimits}
\def\NL{\mathop{N\!L}\nolimits}
\def\CH{\mathop{\mathrm{CH}}\nolimits}
\def\proetale{{pro\text{-}\acute{e}tale}}
\def\etale{{\acute{e}tale}}
\def\QCoh{\mathit{QCoh}}
\def\Ker{\mathop{\mathrm{Ker}}}
\def\Im{\mathop{\mathrm{Im}}}
\def\Coker{\mathop{\mathrm{Coker}}}
\def\Coim{\mathop{\mathrm{Coim}}}

% Boxtimes
%
\DeclareMathSymbol{\boxtimes}{\mathbin}{AMSa}{"02}

%
% Macros for moduli stacks/spaces
%
\def\QCohstack{\mathcal{QC}\!\mathit{oh}}
\def\Cohstack{\mathcal{C}\!\mathit{oh}}
\def\Spacesstack{\mathcal{S}\!\mathit{paces}}
\def\Quotfunctor{\mathrm{Quot}}
\def\Hilbfunctor{\mathrm{Hilb}}
\def\Curvesstack{\mathcal{C}\!\mathit{urves}}
\def\Polarizedstack{\mathcal{P}\!\mathit{olarized}}
\def\Complexesstack{\mathcal{C}\!\mathit{omplexes}}
% \Pic is the operator that assigns to X its picard group, usage \Pic(X)
% \Picardstack_{X/B} denotes the Picard stack of X over B
% \Picardfunctor_{X/B} denotes the Picard functor of X over B
\def\Pic{\mathop{\mathrm{Pic}}\nolimits}
\def\Picardstack{\mathcal{P}\!\mathit{ic}}
\def\Picardfunctor{\mathrm{Pic}}
\def\Deformationcategory{\mathcal{D}\!\mathit{ef}}
\newenvironment{reference}{}{}
\newenvironment{slogan}{}{}
\newenvironment{history}{}{}
\usepackage{cleveref}
\crefname{originallemma}{original lemma}{original lemmas}
\Crefname{originallemma}{Original lemma}{Original lemmas}
\makeatletter
\newlabel{lemma-characterize-geometrically-irreducible}{{038I}{source}{}{sourcepointer-1}{}}
\newlabel{lemma-characterize-geometrically-irreducible@cref}{{[originallemma][0][]038I}{[1][0][]source}{}{}{}}
\makeatother
\crefname{originallemma}{original lemma}{original lemmas}
\Crefname{originallemma}{Original lemma}{Original lemmas}
\makeatletter
\newlabel{lemma-image-irreducible}{{04KX}{source}{}{sourcepointer-2}{}}
\newlabel{lemma-image-irreducible@cref}{{[originallemma][0][]04KX}{[1][0][]source}{}{}{}}
\makeatother
\crefname{originallemma}{original lemma}{original lemmas}
\Crefname{originallemma}{Original lemma}{Original lemmas}
\makeatletter
\newlabel{morphisms-lemma-base-change-finite}{{01WL}{source}{}{sourcepointer-3}{}}
\newlabel{morphisms-lemma-base-change-finite@cref}{{[originallemma][0][]01WL}{[1][0][]source}{}{}{}}
\makeatother
\crefname{originallemma}{original lemma}{original lemmas}
\Crefname{originallemma}{Original lemma}{Original lemmas}
\makeatletter
\newlabel{morphisms-lemma-image-nowhere-dense-finite}{{03HX}{source}{}{sourcepointer-4}{}}
\newlabel{morphisms-lemma-image-nowhere-dense-finite@cref}{{[originallemma][0][]03HX}{[1][0][]source}{}{}{}}
\makeatother
\title{An irreducible scheme with a dense set of rational points is geometrically irreducible}
\author{The Stacks Project Authors}
\date{stacks-project:varieties.tex:local-snapshot-20261011; GNU FDL 1.2 or later}
\begin{document}
\maketitle
\noindent\textit{Editorial note.} All delivered mathematical text is the author's own, copied from the pinned Stacks Project chapter file varieties.tex; no statement or proof step was rewritten, reordered or removed.\par\medskip
\noindent\textit{Editorial note.} The proof is complete as written and ends with the contradiction; no step is referred to the reader.\par\medskip
\noindent\textit{Editorial note.} The counted spans carry 4 navigation stubs and no bibliography citation.\par\medskip
\noindent\textit{Editorial note.} The preamble is the chapter's retained minimal preamble (sources/stacks-varieties/retained-preamble.tex lines 1-77), a verbatim excerpt of the native preamble.\par\medskip
\noindent\textit{Editorial note.} The label lemma-irreducible-dense-rational-points-geometrically-irreducible is defined only in varieties.tex, so no cross-chapter restatement had to be weighed for this card.\par\medskip
\noindent\textit{Allowed standard prerequisites (editorial).}\begin{itemize}
\item Characterization of geometric irreducibility by finite separable field extensions. Same chapter, Lemma lemma-characterize-geometrically-irreducible, invoked at proof lines 1681-1682. Supplies that it is enough to show irreducibility of the base change along finite extensions.
\item Irreducible components of a base change project onto irreducible components. Same chapter, Lemma lemma-image-irreducible, invoked at proof lines 1683-1684. Supplies that both chosen components project onto X.
\item Image of a nowhere dense subset under a finite morphism is nowhere dense, and base change of a finite morphism is finite. Morphisms, Lemmas morphisms-lemma-image-nowhere-dense-finite and morphisms-lemma-base-change-finite, invoked at proof lines 1686-1690. Supplies that the image of the intersection is not dense, so a rational point can be found outside it.
\end{itemize}\medskip
\ifdefined\hypertarget\hypertarget{sourcepointer-1}{}\fi
\noindent\textit{Source reference.} Original lemma 038I: sources/stacks-varieties/pinned-original-source.tex, source lines 1542--1542. Pointer outside the counted span, invoked once for context; the load-bearing argument of the counted statement is the author's proof reproduced in full in the proof span.\par\medskip
\ifdefined\hypertarget\hypertarget{sourcepointer-2}{}\fi
\noindent\textit{Source reference.} Original lemma 04KX: sources/stacks-varieties/pinned-original-source.tex, source lines 1622--1622. Pointer outside the counted span, invoked once for context; the load-bearing argument of the counted statement is the author's proof reproduced in full in the proof span.\par\medskip
\ifdefined\hypertarget\hypertarget{sourcepointer-3}{}\fi
\noindent\textit{Source reference.} Original lemma 01WL: sources/stacks-chow/tags.txt, source lines 2486--2486. Pointer outside the counted span, invoked once for context; the load-bearing argument of the counted statement is the author's proof reproduced in full in the proof span.\par\medskip
\ifdefined\hypertarget\hypertarget{sourcepointer-4}{}\fi
\noindent\textit{Source reference.} Original lemma 03HX: sources/stacks-chow/tags.txt, source lines 4583--4583. Pointer outside the counted span, invoked once for context; the load-bearing argument of the counted statement is the author's proof reproduced in full in the proof span.\par\medskip
\begin{lemma}
\label{lemma-irreducible-dense-rational-points-geometrically-irreducible}
Let $k$ be a field. Let $X$ be a scheme over $k$. If $X$ is irreducible
and has a dense set of $k$-rational points, then $X$ is geometrically
irreducible.
\end{lemma}

\begin{proof}
Let $k'/k$ be a finite extension of fields and let $Z, Z' \subset X_{k'}$
be irreducible components. It suffices to show $Z = Z'$, see
Lemma \ref{lemma-characterize-geometrically-irreducible}.
By Lemma \ref{lemma-image-irreducible} we have $p(Z) = p(Z') = X$
where $p : X_{k'} \to X$ is the projection. If $Z \not = Z'$
then $Z \cap Z'$ is nowhere dense in $X_{k'}$ and hence $p(Z \cap Z')$
is not dense by
Morphisms, Lemma \ref{morphisms-lemma-image-nowhere-dense-finite};
here we also use that $p$ is a finite morphism as the base change
of the finite morphism $\Spec(k') \to \Spec(k)$, see
Morphisms, Lemma \ref{morphisms-lemma-base-change-finite}.
Thus we can pick a $k$-rational point $x \in X$ with
$x \not \in p(Z \cap Z')$. Since the residue field of $x$ is $k$
we see that $p^{-1}(\{x\}) = \{x'\}$ where $x' \in X_{k'}$
is a point whose residue field is $k'$. Since $x \in p(Z) = p(Z')$
we conclude that $x' \in Z \cap Z'$ which is the contradiction
we were looking for.
\end{proof}

\end{document}
